\chapter{複素関数論}
理論の概観を復習していこう。
\ben
\item 複素関数論では\tf{正則関数}が重要です。その定義は単に, \textcolor{blue}{すべての点で複素微分可能である}\footnote{実はこの条件から$C^1$級であることが従います。}というものです。$\ovl{z}$でさえこの条件を満たしませんが, 正則関数からは定義だけでは予想しにくい多くの性質が導かれます。まず, 正則関数の特徴づけとして, \tf{Cauchy-Riemannの関係式}$u_x = v_y, v_x = -u_y$を満たすことが示せます。
\item $\mb{C}$上の線積分が定義される。曲線というものは連続写像$z(t):[a,b]\subset \mb{C}$に他ならないが, 複雑な曲線が登場するわけでもなく, 区分的$C^1$級程度の条件でよい。このような曲線上で, $\sum f(z(t_i)) (z(t_{i+1}) - z(t_i))$というRiemann和的な量の極限として線積分が定義される。この量は複素数であるけれども, そのほかには微積分と変わりのない定義であるため, 微積分と同様の定理は成り立つ(たとえば, \tf{Greenの公式}など)。 しかし複素解析特有の積分定理が多く, いずれも大変重要である。
\item まず\tf{Cauchy積分定理}が示される。これは$D\cup \del D$を含む領域で正則な関数を領域の境界上で積分すると0になるという定理である。これにより, 積分路を変形する技術が可能になり, 簡単な積分が計算可能になったり, 偏角の原理を用いることで代数学の基本定理が証明できたりする。なお, この定理の逆は\tf{Moreraの定理}といい, Moreraの定理を用いることで正則関数の広義一様収束極限が正則であることが示せる。
\item \tf{Cauchyの積分公式}は最重要である。これは正則関数の大域的な性格を引き出す定理である。これにより分かることはいくつもある:  (i)正則関数は解析的, つまり冪級数展開できる (ii) \tf{Goursatの定理}: 正則関数は無限回微分可能 (iii) \tf{Cauchyの係数評価}: 冪級数展開の係数は$a_n = \int_{|\zeta-z|=R} \frac{f(\zeta)}{(\zeta - z)^{n+1}}\frac{d\zeta}{2\pi i n!}$と書け, $a_n \leq R^{-n}\max_{|\zeta - z|=R}|f(\zeta)|$ (iv) \tf{平均値の定理}:$f(c) = \int_{0}^{2\pi} f(c+re^{i\theta})\,\frac{d\theta}{2\pi}$   などの結果が導かれ, (iii)より\tf{Liouvilleの定理}が, (iv)から積分の絶対値を見ると\tf{最大値の原理}が示される。そもそも, 解析性と$C^{\infty}$級が同じになること自体, 実解析では起こらないことであることにも注意したい。
\item これまでは正則関数の性質を見てきたが, この次に孤立特異点$z=c$を持つ\footnote{$z=c$で正則な場合も点$c$を孤立特異点と呼ぶ。}関数についても考える。このような関数は$\frac{\sin{z}}{z}, \frac{1}{z(z^2-1)}, e^{\frac{1}{z}}$などがあり, いずれも0が孤立特異点であるが, それぞれ\tf{除去可能特異点}, \tf{極}, \tf{真性特異点}と呼ばれる特異点である。しかし共通に言えることが一つある。今, 0が孤立特異点である場合を考えると, このような孤立特異点の持つ関数関数は, それが$0<|z|<R$で正則であるのなら,  \tf{$0<|z|<R$で広義一様収束するLaurent級数に展開可能}である\footnote{このときの係数は$f(z)$にしか依らない}。"広義"とあるから,これを示すときには円環領域$0<r_1\leq |z|\leq r_2 < R$で考えなければならない。この任意の$r_1,r_2$について考えるというところに孤立性が効いていることに注意。与えられた関数のLaurent展開の$z^m$の係数$a_m$が知りたければ, Cauchy積分公式から$a_m = \int_{|z| = r} f(z)z^{-(m+1)}\frac{dz}{2\pi i}$を計算すればよい。$z=c$が孤立特異点のときも同様である。覚えてもらいたいことは, この広義一様収束性である。
\item 3種類の特異点があったが, これらの正確な定義は, その点のまわりのLaurent展開の負冪の係数がどうなっているかで分けられている。負冪がないとき,負冪が有限個だけあるとき, 負冪が無限個あるとき, それぞれ 除去可能特異点, 極, 真性特異点というのである。また, これらには計算面で有用な特徴付けがあり, 実際以下のようになっている: (i) \tf{Riemannの除去可能定理}: $z=c$が除去可能特異点$\iff \lim_{z\to c} f(z)$が有限確定\footnote{もしくは適当な$0<|z-c|<r$で$|f(z)|$が有界(その他にもいろいろある)} 　(ii) $z=c$が極$\iff$ $\lim_{z\to c}|f(z)| = \infty$ (iii) \tf{Casorati-Weierstrassの定理}:$z=c$が真性特異点$\iff$ $c$の任意の近傍の$f$による像は$\mb{C}$上稠密 ($\iff$ いかなる$\alpha\in \mb{C}$に対しても, ある点列$z_n\to c$が存在して$f(z_n)\to \alpha$)\footnote{$(\spadesuit\spadesuit)$ 実際のところ, この近傍の像は$\mb{C}-\{ 1点\}$であるか$\mb{C}$であることが知られている。つまり, 真性特異点のまわりで取らない複素数があるとすれば, それはたったひとつしかない。これは\tf{Picardの定理}で有名である。 }。
\item とくに$a_{-1}$の係数を\tf{$c$における留数}といい, $a_{-1} = \mr{Res}_{z=c}f(z)dz$と表す。これを調べることで多くの積分を計算できるようになるというのが\tf{留数定理}である。正確な主張は次の通りである: \tf{$f(z)$は単純閉曲線$C$の内部に孤立特異点$c_1,\cdots, c_N$をもつほかは, $C$の内部と周をこめて正則とするとき, $\int_{C}f(z)dz = 2\pi i\sum_{j=1}^{N} \mr{Res}_{z=c_j} f(z)dz$}。留数定理は, $\int_{C}f(z)dz$を求めるための一連の計算を簡略化する定理である。その手続きは次の4段階からなる: (i) $c_j$の周りの十分小さい円$C_j$を描く (ii)積分路変形によって$\int_{C}f(z)dz = \sum_{j=1}^N \int_{C_j} f(z)dz$ (iii)$z=c_j$のLaurent展開を記録し, \ao{広義一様収束によって項別積分する} (iv) $\int_{|z|=r} z^n\, dz$は$n=-1$,$n\neq -1$でそれぞれ$2\pi i,0$だから, $\int_{C_j}f(z)dz = 2\pi i \mr{Res}_{z=c_j} f(z)dz$
\item \tf{積分路変形原理}と, \tf{留数定理}と, \tf{Cauchyの係数評価}と, \tf{Jordanの不等式}($x\leq \frac{\pi}{2}\sin{x}$ on $[0,\frac{\pi}{2}]$) でかなり多くの積分が計算できる。後半2つは広義積分に関連している。\tf{積分路変形原理}を使えば, 特異点を跨がない限り経路を変更できる。\tf{留数定理}を使えば, 経路内部の留数計算をするだけで閉曲線の積分計算ができる。\tf{Cauchyの係数評価}や\tf{Jordanの不等式}を使えば, 無限遠点の周りを一周する経路の積分が$0$に収束することがしばしば言える。経路としては, 上半円, 円環の上半分, 長方形, 扇形などがある。分枝を取るものに関しては, その分枝が1価に定義できる領域を選択しなければならない。
\een

\newpage
